\documentclass[11pt]{article}
\usepackage[margin=1in]{geometry}
\usepackage{graphicx}
\usepackage{parskip}
\usepackage{titlesec}
\usepackage{xcolor}

\titleformat{\section}{\large\bfseries}{}{0pt}{}
\titlespacing*{\section}{0pt}{12pt}{4pt}

\setlength{\parindent}{0pt}

\begin{document}

{\LARGE \bfseries Tracing Low-Confidence Edges: False Positive Analysis}

\vspace{4pt}
Our connectivity estimates include false positives -- pairs that look
connected but probably aren't. Rather than trying to label each edge simply
``real'' or ``fake,'' I've been looking for signals that shift our
\emph{confidence} one way or the other. None of these are hard rules: we are
not certain a real edge is never bidirectional, or that a fake edge always
is -- there are exceptions on both sides. What follows is a set of
percentage-based tendencies, not a certainty test.

\section*{1. Timing}
\includegraphics[width=0.95\textwidth]{../../figures/fig_R8b_signal.png}

A real synapse has a fixed delay ($\sim$1.5\,ms here); a shared-input fake
should look closer to instantaneous. This compares average connection
strength across time lags for likely-real vs.\ likely-false edges, on raw
spikes and on the calcium signal.
\emph{How to read it:} on spikes, the two curves separate cleanly. Through
calcium they collapse onto each other -- the calcium sensor blurs time by
about 3\,ms, coarser than the 1.5\,ms delay, so this signal isn't usable on
calcium data. A limit of the measurement, not the idea.
\emph{Computed by:} a lagged coupling estimate at 10 lags (0.1--6\,ms) from
cached moments, averaged by class at each lag.

\section*{2. Timing, as a plain count}
\includegraphics[width=0.95\textwidth]{../../figures/fig_R8b_signal_counts.png}

Same data, simpler question: for each pair, at \emph{which} lag is its
coupling strongest? I counted what fraction of each group peaks at each lag.
\emph{How to read it:} on spikes, nearly all likely-real edges peak exactly
at the true delay. Through calcium, both groups peak at the same lags and
overlap. Easier to read than the averaged curve above; same conclusion.

\section*{3. Directionality}
\includegraphics[width=0.95\textwidth]{../../figures/fig_R8_violin_n12500.png}

A different candidate signal: a real synapse mostly goes one direction;
shared input tends to make \emph{both} directions look connected. I'm not
assuming this is a hard rule -- some real connections genuinely go both
ways -- but I checked whether it shifts the odds.
\emph{How to read it:} each violin is the distribution of a directionality
score (0 = both directions equally strong, 1 = one direction only).
Likely-real edges sit high; likely-false ones sit noticeably lower -- not a
clean split, but a real shift.
\emph{Computed by:} for each pair, $|A[i,j]-A[j,i]| / (|A[i,j]|+|A[j,i]|)$.

\section*{4. Directionality, as a plain percentage}
\includegraphics[width=0.8\textwidth]{../../figures/fig_R8_counts_n1250_local.png}

Instead of the full distribution, the plain question: is the reverse
direction \emph{also} flagged as connected? Likely-real edges are
``bidirectional'' about 10\% of the time; likely-false ones, about 38\%.
Again -- a shift in odds, not a certainty. Some real edges will still fall
in that 38\%, and some fakes in the 10\%.

\section*{5. Where the uncertainty comes from}
\includegraphics[width=0.95\textwidth]{../../figures/fig_way2_motifs_n1250.png}

To understand this better: does a false positive's strength depend on how
many other neurons feed into both members of a pair -- and does it matter
whether we can actually \emph{record} those driver neurons?
\emph{How to read it:} three views of the same false positives -- direct
shared drivers (left), indirect two-hop shared drivers (middle), driver
\emph{strength} instead of headcount (right). With a plain headcount, hidden
drivers matter slightly more. Once weighted by actual driver strength, that
gap nearly disappears -- the confound tracks total exposure to shared input,
not whether we happened to record its source.

\section*{Where this leaves us}
None of these signals cleanly separate real from fake on their own. What
they give us is a set of odds-shifting features -- timing (spikes only),
directionality (works throughout, including calcium) -- plus a clearer
picture of what drives the confound in the first place. Next step: combine
these into a single per-edge confidence score rather than a hard threshold.

\end{document}
